% Part: lambda-calculus
% Chapter: syntax
% Section: alpha

\documentclass[../../../include/open-logic-section]{subfiles}

\begin{document}

\olfileid{lam}{syn}{alp}
\olsection{تحويل~$\alpha$}

ما العلاقة بين~$\lambd[x][x]$ و$\lambd[y][y]$؟ كلاهما يمثل دالة
الهوية، مع أنهما بالطبع حدان مختلفان من الوجهة التركيبية. ولا يختلفان إلا
في اسم المتغير المقيد، وأحدهما ناتج ``إعادة تسمية'' المتغير المقيد في
الآخر. وهذا ما يسمى \emph{تحويل~$\alpha$}.

\begin{defn}[تغيير المتغير المقيد، $\aconvone$]
  إذا احتوى حد~$M$ واقعةً لـ$\lambd[x][N]$، وكان~$y \notin
  \FV{N}$، وكان~$\Subst{N}{y}{x}$ معرّفًا، فإن استبدال هذه الواقعة بـ
  \begin{equation*}
    \lambd[y][\Subst{N}{y}{x}]
  \end{equation*}
  بحيث ينتج~$M'$ يسمى \emph{تغييرًا للمتغير المقيد}، ويكتب
  $M \redone[\alpha] M'$.
\end{defn}

\begin{defn}[توافق العلاقة]
  يقال إن علاقة~$R$ على الحدود \emph{متوافقة} إذا استوفت الشروط الآتية:
  \begin{enumerate}
  \item إذا كان~$R N N'$ فإن~$R \lambd[x][N] \lambd[x][N']$
  \item إذا كان~$R P P'$ فإن~$R (PQ) (P'Q)$
  \item إذا كان~$R Q Q'$ فإن~$R (PQ) (PQ')$
  \end{enumerate}
\end{defn}

فلنعد الآن صياغة التعريف:
\begin{defn}[تغيير المتغير المقيد، $\aconvone$]
  \emph{تغيير المتغير المقيد} ($\redone[\alpha]$) هو أصغر علاقة متوافقة
  على الحدود تستوفي الشرط الآتي:
  \begin{align*}
    &\lambd[x][N] \redone[\alpha] \lambd[y][\Subst{N}{y}{x}] && \text{إذا
      كان~$x \neq y$ و$y \notin \FV{N}$} \\
    & &&\text{وكان~$\Subst{N}{y}{x}$ معرّفًا}
  \end{align*}
\end{defn}

تعني ``أصغر'' هنا أن العلاقة لا تحتوي إلا الأزواج التي يقتضيها التوافق
والشرط الإضافي، ولا تحتوي سواها. ولذلك يمكن تعريف هذه العلاقة أيضًا كما
يأتي:

\begin{defn}[تغيير المتغير المقيد، $\aconvone$] \ollabel{defn:aconvone}
  يعرّف \emph{تغيير المتغير المقيد} ($\aconvone$) استقرائيًا كما يأتي:
  \begin{enumerate}
  \item إذا كان~$N \aconvone N'$ فإن~$\lambd[x][N] \aconvone
    \lambd[x][N']$ \ollabel{defn:aconvone1}
  \item إذا كان~$P \aconvone P'$ فإن~$(PQ) \aconvone (P'Q)$ \ollabel{defn:aconvone2}
  \item إذا كان~$Q \aconvone Q'$ فإن~$(PQ) \aconvone (PQ')$ \ollabel{defn:aconvone3}
  \item إذا كان~$x \neq y$ و$y \notin \FV{N}$ وكان~$\Subst{N}{y}{x}$
    معرّفًا، فإن~$\lambd[x][N] \redone[\alpha]
    \lambd[y][\Subst{N}{y}{x}]$.
    \ollabel{defn:aconvone4}
  \end{enumerate}
\end{defn}

التعريفات متكافئة، لكننا نترك البرهان تمرينًا. وسنستعمل من الآن فصاعدًا
التعريف الاستقرائي.

\begin{defn}[تحويل~$\alpha$، $\aconv$]
  \emph{تحويل~$\alpha$} ($\aconv$) هو أصغر علاقة انعكاسية ومتعدية على
  الحدود تحتوي~$\aconvone$.
\end{defn}

وكما سبق، تعني ``أصغر'' أن العلاقة لا تحتوي إلا الأزواج التي يقتضيها
التعدي و$\aconvone$، وهذا يقود إلى التعريف المكافئ الآتي:

\begin{defn}[تحويل~$\alpha$، $\aconv$] \ollabel{defn:aconv}
  يعرّف \emph{تحويل~$\alpha$} ($\aconv$) استقرائيًا كما يأتي:
  \begin{enumerate}
  \item إذا كان~$P \aconv Q$ و$Q \aconv R$، فإن~$P \aconv R$.
    \ollabel{defn:aconv1}
  \item إذا كان~$P \aconvone Q$، فإن~$P \aconv Q$. \ollabel{defn:aconv2}
  \item $P \aconv P$. \ollabel{defn:aconv3}
  \end{enumerate}
\end{defn}

\begin{ex}
  يتحول~$\lambd[x][f x]$ بتحويل~$\alpha$ إلى~$\lambd[y][f
  y]$، والعكس صحيح. وبعبارة غير رسمية، كلاهما دالة تقبل وسيطًا وتعيد
  نتيجة تطبيق~$f$ على ذلك الوسيط، مع الإحالة إلى متغير البيئة~$f$.

  ولا يتحول~$\lambd[x][f x]$ بتحويل~$\alpha$ إلى~$\lambd[x][g
    x]$. وبعبارة غير رسمية، يحيل الحدان إلى متغيري البيئة~$f$ و$g$ على
  الترتيب، وهذا يجعلهما دالتين مختلفتين: فسلوكهما يختلف في البيئات التي
  تختلف فيها~$f$ و$g$.
\end{ex}

\begin{prob}
  هل كل زوج من أزواج الحدود الآتية قابل لتحويل~$\alpha$؟
  \begin{enumerate}
  \item $\lambd[x][\lambd[y][x]]$ و$\lambd[y][\lambd[x][y]]$
  \item $\lambd[x][\lambd[y][x]]$ و$\lambd[c][\lambd[b][a]]$
  \item $\lambd[x][\lambd[y][x]]$ و$\lambd[c][\lambd[b][a]]$
  \end{enumerate}
\end{prob}

\begin{lem}\ollabel{lem:fv-one}
  إذا كان~$P \aconvone Q$ فإن~$\FV{P} = \FV{Q}$.
\end{lem}

\begin{proof}
  بالاستقراء على~!!{derivation} لـ$P \aconvone Q$.
  \begin{enumerate}
  \item إذا كانت القاعدة الأخيرة هي~\olref{defn:aconvone4}، كان~$P$ من
    الصورة~$\lambd[x][N]$ و$Q$ من الصورة
    $\lambd[y][\Subst{N}{y}{x}]$، حيث~$x \neq y$ و$y \notin \FV{N}$
    و$\Subst{N}{y}{x}$ معرّف. ونميز حالتين بحسب كون~$x \in \FV{N}$:
    \begin{enumerate}
    \item إذا كان~$x \in \FV{N}$، فعندئذ:
      \begin{align*}
        \FV{\lambd[y][\Subst{N}{y}{x}]} &=
        \FV{\Subst{N}{y}{x}} \setminus \{y\} \\
        & = ((\FV{N} \setminus \{x\}) \cup \{y\}) \setminus \{y\}
         && \text{ بموجب \olref[sub]{thm:infv}} \\
        & = \FV{N} \setminus \{x\} \\
        & = \FV{\lambd[x][N]}
      \end{align*}
    \item إذا كان~$x \notin \FV{N}$، فعندئذ:
      \begin{align*}
        \FV{\lambd[y][\Subst{N}{y}{x}]}
        & = \FV{\Subst{N}{y}{x}} \setminus \{y\} \\
        & = \FV{N} \setminus \{x\}
         && \text{بموجب \olref[sub]{thm:notinfv}} \\
        & = \FV{\lambd[x][N]}.
      \end{align*}
    \end{enumerate}
  \item تترك الحالات الثلاث الأخرى تمارين. 
  \end{enumerate}
\end{proof}

\begin{prob}
  أكمل برهان~\olref[lam][syn][alp]{lem:fv-one}.
\end{prob}

\begin{lem}\ollabel{lem:inv}
  إذا كان~$P \aconvone Q$ فإن~$Q \aconvone P$.
\end{lem}

\begin{proof}
  بالاستقراء على~!!{derivation} لـ$P \aconvone Q$.
  \begin{enumerate}
  \item إذا كانت القاعدة الأخيرة هي~\olref{defn:aconvone4}، كان~$P$ من
    الصورة~$\lambd[x][N]$ و$Q$ من الصورة
    $\lambd[y][\Subst{N}{y}{x}]$، حيث~$x \neq y$ و$y \notin \FV{N}$
    و$\Subst{N}{y}{x}$ معرّف. أولًا، لدينا~$y \notin
    \FV{\Subst{N}{y}{x}}$ بموجب~\olref[sub]{thm:clr}. وبموجب
    \olref[sub]{thm:inv} نعلم أن~$\Subst{\Subst{N}{y}{x}}{x}{y}$ ليس
    معرّفًا فحسب، بل يساوي~$N$ أيضًا. ثم نحصل بموجب
    \olref{defn:aconvone4} على~$\lambd[y][\Subst{N}{y}{x}]
    \aconvone \lambd[x][\Subst{\Subst{N}{y}{x}}{x}{y}] =
    \lambd[x][N]$.
  \end{enumerate}
\end{proof}

\begin{prob}
  أكمل برهان~\olref[lam][syn][alp]{lem:inv}
\end{prob}

\begin{thm}
  تحويل~$\alpha$ علاقة تكافؤ على الحدود، أي إنه انعكاسي ومتناظر ومتعدٍ.
\end{thm}

\begin{proof}
  \begin{enumerate}
  \item يمكن تغيير كل حد~$M$ إلى~$M$ بإجراء \emph{صفر} من تغييرات
    المتغيرات المقيدة.
  \item إذا تحول~$P$ بتحويل~$\alpha$ إلى~$Q$ بسلسلة من تغييرات
    المتغيرات المقيدة، أمكننا انطلاقًا من~$Q$ أن نعكس هذه التغييرات
    ببساطة (بموجب~\olref{lem:inv}) وبترتيب معاكس، فنحصل على~$P$.
  \item إذا تحول~$P$ بتحويل~$\alpha$ إلى~$Q$ بسلسلة من تغييرات
    المتغيرات المقيدة، وتحول~$Q$ إلى~$R$ بسلسلة أخرى، أمكننا تغيير~$P$
    إلى~$R$ بتطبيق السلسلة الأولى ثم الثانية.
  \end{enumerate}
\end{proof}

ومن الآن فصاعدًا نقول إن~$M$ و$N$ \emph{متكافئان بحسب~$\alpha$}، ونكتب
$M \aeq N$، إذا وفقط إذا تحول~$M$ بتحويل~$\alpha$ إلى~$N$ (وهو، كما
برهنا للتو، حاصل إذا وفقط إذا تحول~$N$ بتحويل~$\alpha$ إلى~$M$).

\begin{thm}\ollabel{thm:fv}
  إذا كان~$M \aeq N$، فإن~$\FV{M} = \FV{N}$.
\end{thm}

\begin{proof}
  ينتج مباشرة من~\olref{lem:fv-one}.
\end{proof}

\begin{lem}\ollabel{lem:sub:R}
  إذا كان~$R \aeq R'$ وكان~$\Subst{M}{R}{y}$ معرّفًا، فإن
  $\Subst{M}{R'}{y}$ معرّف ومكافئ بحسب~$\alpha$ لـ$\Subst{M}{R}{y}$.
\end{lem}

\begin{proof}
  تمرين.
\end{proof}

\begin{prob}
  أثبت~\olref[lam][syn][alp]{lem:sub:R}.
\end{prob}

تذكّر أن الإحلال في~\olref[sub]{sec} غير معرّف في بعض الحالات؛ لكننا
نستطيع، باستعمال تحويل~$\alpha$ على الحدود، أن نجعل الإحلال معرّفًا
دائمًا بإعادة تسمية المتغيرات المقيدة. وتحفظ النتيجة التكافؤ بحسب
$\alpha$، كما تبين المبرهنة الآتية:

\begin{thm}\ollabel{thm:sub}
  لكل~$M$ و$R$ و$y$ يوجد~$M'$ بحيث~$M \aeq M'$ ويكون
  $\Subst{M'}{R}{y}$ معرّفًا. وفضلًا على ذلك، إذا وجد زوج آخر يحقق
  $M'' \aeq M$ و$R''$ وكان~$\Subst{M''}{R''}{y}$ معرّفًا و$R'' \aeq
  R$، فإن~$\Subst{M'}{R}{y} \aeq \Subst{M''}{R''}{y}$.
\end{thm}

\begin{proof}
  بالاستقراء على تكوين~$M$:
  \begin{enumerate}
  \item إذا كان~$M$ متغيرًا~$z$: تمرين.
  \item افترض أن~$M$ من الصورة~$\lambd[x][N]$. اختر متغيرًا~$z$
    يختلف عن~$x$ و$y$ ويحقق~$z \notin \FV{N}$ و$z \notin \FV{R}$.
    وبفرض الاستقراء يوجد~$N'$ بحيث~$N' \aeq N$ ويكون
    $\Subst{N'}{z}{x}$ معرّفًا. وعندئذ أيضًا~$\lambd[x][N] \aeq
    \lambd[x][N']$ بموجب
    \olref{defn:aconvone}\olref{defn:aconvone1}. والآن~$\lambd[x][N']
    \aeq \lambd[z][\Subst{N'}{z}{x}]$ بموجب
    \olref{defn:aconvone}\olref{defn:aconvone4}. ويمكننا فعل ذلك لأن
    $z \ne x$ و$z \notin \FV{N'}$ و$\Subst{N'}{z}{x}$ معرّف. وأخيرًا،
    يكون~$\Subst{\lambd[z][\Subst{N'}{z}{x}]}{R}{y}$ معرّفًا لأن~$z
    \neq y$ و$z \notin \FV{R}$.

    وفضلًا على ذلك، إذا وجد~$N''$ و$R''$ آخران يستوفيان الشروط نفسها،
    كان
    \begin{multline*}
      \Subst{(\lambd[z][\Subst{N''}{z}{x}])}{R''}{y} =\\
    \begin{aligned}
      &= \lambd[z][\Subst{\Subst{N''}{z}{x}}{R''}{y}] \\
      &= \lambd[z][\Subst{\Subst{N''}{z}{x}}{R}{y}] && \text{بموجب
        \olref{lem:sub:R}}\\
      &=\lambd[z][\Subst{\Subst{N'}{z}{x}}{R}{y}]
      && \text{بفرض الاستقراء}\\
      &=\Subst{(\lambd[z][\Subst{N'}{z}{x}])}{R}{y}
    \end{aligned}
    \end{multline*}
    \item إذا كان~$M$ من الصورة~$(PQ)$: تمرين.
  \end{enumerate}
\end{proof}

\begin{prob}
  أكمل برهان~\olref[lam][syn][alp]{thm:sub}.
\end{prob}

\begin{cor}\ollabel{cor:sub}
  لكل~$M$ و$R$ و$y$ يوجد زوج~$M'$ و$R'$ بحيث~$M' \aeq M$ و$R \aeq
  R'$ ويكون~$\Subst{M'}{R'}{y}$ معرّفًا. وفضلًا على ذلك، إذا وجد زوج
  آخر يحقق~$M'' \aeq M$ و$R''$ وكان~$\Subst{M''}{R''}{y}$ معرّفًا،
  فإن~$\Subst{M'}{R'}{y} \aeq \Subst{M''}{R''}{y}$.
\end{cor}

\begin{proof}
  ينتج مباشرة من~\olref{thm:sub}.
\end{proof}

\end{document}
